
Batch Normalization amplifies worst-group accuracy gaps by 9.41 percentage points compared to Layer Normalization at matched average accuracy on synthetic spurious correlation tasks (p = 5.$43\times 10^{-5}$, Cohen's d = 3.94). This effect operates via a gradient flow mechanism: Batch Normalization shows 26\% higher gradient magnitude toward spurious-aligned samples during early training (epochs 0--19) compared to Layer Normalization. These findings are based on proof-of-concept experiments using synthetic data with 90\% spurious correlation; real dataset validation (Waterbirds, CelebA) is pending. The results suggest that normalization layer choice affects fairness metrics in models trained on datasets with spurious correlations. Layer Normalization's instance-level normalization eliminates batch-level spurious signal aggregation present in Batch Normalization. This study introduces accuracy-matched temporal comparison as a method for isolating architectural effects during training.
